\clearpage
\thispagestyle{plain}
\begin{center}
{\fontsize{14}{18}\selectfont\bfseries Abstract\par}
\end{center}

\vspace{0.1cm}

\setstretch{1.5}
\setlength{\parskip}{3pt}
This dissertation investigates the value relevance of financial accounting information—specifically earnings per share (EPS) and book value per share (BVPS)—among prominent Nepalese commercial banks, and evaluates whether these valuation relationships underwent a structural regime shift following the commercial banking consolidation wave of early 2023. Grounded in the residual income valuation (RIV) framework of Ohlson (1995) and clean-surplus accounting theory, the study resolves critical methodological limitations in prior Nepalese literature by eliminating circular price-to-earnings (P/E) regressors, removing look-ahead bias via post-reporting-deadline prices under Section 59(2) of BAFIA ($MPS^{RD}$), and correcting for small-cluster downward bias. A longitudinal panel of five commercial banking series is examined across six fiscal years (FY 2076/77–2081/82, $N = 28$ clean bank-years under a 5-day freshness rule). Econometric estimation employs two-way fixed effects panel regressions with MacKinnon-White HC3 robust standard errors evaluated on Student's $t$ residual degrees of freedom.

In the baseline specification (Model 1), both earnings ($\hat{\beta}_E = 2.095, p = 0.616$) and book value ($\hat{\beta}_B = 0.219, p = 0.908$) exhibit positive but statistically insignificant slopes ($R^2 = 0.948$), showing that common fiscal-year fixed effects absorb macro-temporal price movements. In the regime shift specification (Model 2), pre-consolidation pricing was driven primarily by earnings on point estimates ($\hat{\beta}_E = 5.577, p = 0.140$), but post-consolidation interactions compressed implied slopes to $2.453$ for EPS and $-1.658$ for BVPS. The baseline joint Wald test fails to reject structural stability at the 5\% level ($F(2, 14) = 2.954, p = 0.085$). However, purging the entity-timing asymmetry in Specification R6 ($N = 23$) yields a statistically significant joint test at 5\% ($F(2, 10) = 5.82, p = 0.021$). The evidence demonstrates that the empirical regime shift is sensitivity-dependent rather than conclusive in the baseline, with post-merger credit provisioning and capital dilution coinciding with compressed multiples rather than book-value substitution. The study provides targeted policy implications for regulatory asset quality reviews, trading halt limits, and centralized electronic corporate filings.

\vspace{0.2cm}

\noindent
\textbf{Keywords:} Value Relevance, Commercial Banks, Bank Consolidation, Ohlson Model, Two-Way Fixed Effects.

